% ============================================
% DESEMPENHO DO SSB
% Detecção coerente, como no DSB-SC. O que muda é a banda do filtro de
% recepção e a parcela em quadratura do sinal transmitido.
% ============================================

% Receptor coerente com o filtro de recepção como argumento.
% #1: texto do filtro; #2: ponto destacado, 1 ou 2.
\newcommand{\diagCoe}[2]{%
\begin{tikzpicture}[font=\scriptsize,
  blk/.style={rectangle, draw, rounded corners, minimum height=0.6cm, minimum width=1.2cm, align=center},
  sm/.style={circle, draw, inner sep=0.5pt}]
\node (in) {$s(t)$};
\node[sm, right=0.5cm of in] (sum) {$+$};
\node[above=0.25cm of sum, cRed] (n) {$n(t)$};
\node[blk, right=0.5cm of sum] (bpf) {#1};
\node[sm, right=1.3cm of bpf] (mix) {$\times$};
\node[above=0.25cm of mix] (lo) {$\cos(2\pi f_c t)$};
\node[blk, right=0.5cm of mix] (lpf) {PB $W$};
\node[right=1.3cm of lpf] (out) {$y(t)$};
\draw[->] (in) -- (sum);
\draw[->, cRed] (n) -- (sum);
\draw[->] (sum) -- (bpf);
\draw[->] (bpf) -- (mix);
\draw[->] (lo) -- (mix);
\draw[->] (mix) -- (lpf);
\draw[->] (lpf) -- (out);
\coordinate (p1) at ($(bpf.east)!0.5!(mix.west)$);
\coordinate (p2) at ($(lpf.east)!0.5!(out.west)$);
\ifnum#2=1
  \node[marca, fill=cBlue, text=white] at (p1) {1};
\else
  \node[marca, draw=gray!60, fill=white, text=gray] at (p1) {1};
\fi
\ifnum#2=2
  \node[marca, fill=cGreen, text=white] at (p2) {2};
\else
  \node[marca, draw=gray!60, fill=white, text=gray] at (p2) {2};
\fi
\end{tikzpicture}}

\newcommand{\topoSSB}[1]{{\centering\diagCoe{PF $W$, de\\ $f_c$ a $f_c+W$}{#1}\par}\vspace{2pt}}

\subsection{SSB}

\begin{frame}{Receptor SSB}

\small\compacto
Banda lateral superior, com $m_h$ a transformada de Hilbert de $m$:
\[
s(t) = A_c\big[m(t)\cos(2\pi f_c t) - \tikzmarknode{mh}{m_h(t)}\sin(2\pi f_c t)\big]
\]
\par\vspace{0.45cm}%
\anotacoes{\only<2->{
  \node[nota=cOrange, anchor=north east] (a) at ([xshift=-1.2cm,yshift=-0.15cm]mh.south west) {$M_h(f) = -j\,\mathrm{sgn}(f)\,M(f)$};
  \draw[setinha, cOrange] (a.east) -- (mh.south west);
}}%

\vspace{-0.3cm}
\begin{center}
\begin{tikzpicture}
\begin{axis}[onda, width=9cm, height=1.1cm, xmin=-4, xmax=4, ymin=0, ymax=1.3,
             xlabel={$f$}, ylabel={espectro de $s(t)$}]
  \only<3->{
    \addplot[cBlue, thick, fill=cBlue!20, samples=2] coordinates {(2,0) (2,0.8) (3.2,0)};
    \addplot[cBlue, thick, fill=cBlue!20, samples=2] coordinates {(-2,0) (-2,0.8) (-3.2,0)};
    \node[font=\tiny, anchor=north] at (axis cs:2,0) {$f_c$};
    \node[font=\tiny, anchor=north] at (axis cs:3.2,0) {$f_c+W$};
    \node[font=\tiny, anchor=north] at (axis cs:-2,0) {$-f_c$};
  }
  \only<4->{
    \draw[cGreen, thick, dashed] (axis cs:2,0) -- (axis cs:2,1.1) -- (axis cs:3.2,1.1) -- (axis cs:3.2,0);
    \draw[cGreen, thick, dashed] (axis cs:-2,0) -- (axis cs:-2,1.1) -- (axis cs:-3.2,1.1) -- (axis cs:-3.2,0);
    \node[cGreen, font=\tiny, anchor=south] at (axis cs:2.6,1.1) {filtro de recepção};
  }
\end{axis}
\end{tikzpicture}
\end{center}

\only<1-2>{O sinal ocupa só a banda lateral superior.}%
\only<3>{Banda $W$, metade do DSB-SC.}%
\only<4>{O filtro de recepção acompanha o sinal: largura $W$, de $f_c$ a $f_c+W$.}%
\only<5->{{\centering\diagCoe{PF $W$, de\\ $f_c$ a $f_c+W$}{0}\par}}

\end{frame}

% ============================================

\begin{frame}{SSB, ponto 1: potência do sinal}

\small\compacto
\[
S_i = A_c^2\Big[\overline{m^2}\,\tikzmarknode{c2}{\overline{\cos^2}} + \overline{m_h^2}\,\overline{\sin^2} - 2\,\tikzmarknode{cs}{\overline{m\,m_h\cos\sin}}\Big]
\only<3->{= \frac{A_c^2}{2}\big(\overline{m^2} + \tikzmarknode{mh2}{\overline{m_h^2}}\big)}
\]
\par\vspace{0.7cm}%
\anotacoes{
  \only<2->{
    \node[nota=cBlue, anchor=north east] (a) at ([xshift=-0.2cm,yshift=-0.3cm]c2.south) {$\overline{\cos^2} = \overline{\sin^2} = \tfrac12$};
    \draw[setinha, cBlue] (a.north) -- (c2.south);
    \node[nota=cRed, anchor=north] (b) at ([xshift=0.4cm,yshift=-0.3cm]cs.south) {$\cos\sin = \tfrac12\sin 4\pi f_c t$: média nula};
    \draw[setinha, cRed] (b.north) -- (cs.south);
  }
}%
\onslide<4->{Potência de $m_h$ pela PSD, com $|{-j}\,\mathrm{sgn}(f)| = 1$:}
\onslide<4->{
\[
\overline{m_h^2} = \int_{-\infty}^{\infty} S_m(f)\,\tikzmarknode{hh}{\big|{-j}\,\mathrm{sgn}(f)\big|^2}\,df
\only<5->{= \int_{-\infty}^{\infty} S_m(f)\,df = \overline{m^2}}
\]}%
\par\vspace{0.5cm}%
\anotacoes{\only<5->{
  \node[nota=cOrange, anchor=north] (c) at ([yshift=-0.25cm]hh.south) {Hilbert só muda a fase};
  \draw[setinha, cOrange] (c.north) -- (hh.south);
}}%
\onslide<6->{
\[
\boxed{\,S_i = A_c^2\,\overline{m^2}\,}
\]}

\end{frame}

% ============================================

\begin{frame}{SSB, ponto 1: ruído e SNR}

\topoSSB{1}

\begin{columns}[T]
\column{0.5\textwidth}
\begin{tikzpicture}
\begin{axis}[onda, width=5.6cm, height=1.4cm, xmin=-4, xmax=4, ymin=0, ymax=1.3,
             xlabel={$f$}, ylabel={$S_n(f)$ no ponto 1}]
  \only<1>{\addplot[cRed, thick, domain=-4:4, samples=2] {0.6};}
  \only<2->{
    \addplot[cRed, thick, fill=cRed!20, samples=2] coordinates {(2,0) (2,0.6) (3.2,0.6) (3.2,0)};
    \addplot[cRed, thick, fill=cRed!20, samples=2] coordinates {(-3.2,0) (-3.2,0.6) (-2,0.6) (-2,0)};
    \node[font=\tiny, anchor=north] at (axis cs:2,0) {$f_c$};
    \node[font=\tiny, anchor=north] at (axis cs:-2,0) {$-f_c$};
    \draw[cBlue, <->] (axis cs:2,0.85) -- node[above, font=\tiny] {$W$} (axis cs:3.2,0.85);
  }
  \node[cRed, font=\tiny, anchor=south] at (axis cs:0,0.6) {$N_0/2$};
\end{axis}
\end{tikzpicture}

\vspace{0.2cm}\small
\only<1>{Ruído branco antes do filtro.}%
\only<2->{Duas faixas de largura $W$, metade das do DSB-SC.}

\column{0.5\textwidth}
\small\compacto
\onslide<3->{
\[
\overline{n^2} = 2\times\frac{N_0}{2}\times W = N_0W
\]}%
\onslide<4->{
\[
\left(\frac{S}{N}\right)_1 = \frac{S_i}{N_0W}
\only<5->{= \gamma}
\]}%
\onslide<6->{A banda menor deixa entrar a metade do ruído do DSB-SC: no ponto 1 o SSB já está com $\gamma$, 3\,dB acima do DSB-SC.}
\end{columns}

\end{frame}

% ============================================

\begin{frame}{SSB: PSD do ruído $n_c$}

\begin{columns}[T]
\column{0.55\textwidth}
\begin{tikzpicture}
\begin{axis}[onda, width=5.4cm, height=0.6cm, xmin=-6, xmax=6, ymin=0, ymax=1.3,
             every axis y label/.style={at={(rel axis cs:0,0.3)}, anchor=east, font=\scriptsize},
             xlabel={$f$}, ylabel={$S_n(f-f_c)$}]
  \addplot[cOrange, thick, fill=cOrange!20, samples=2] coordinates {(0,0) (0,0.5) (1,0.5) (1,0)};
  \addplot[cOrange, thick, fill=cOrange!20, samples=2] coordinates {(4.5,0) (4.5,0.5) (5.5,0.5) (5.5,0)};
  \node[font=\tiny, anchor=north] at (axis cs:4.5,0) {$2f_c$};
  \node[font=\tiny, anchor=north] at (axis cs:1,0) {$W$};
\end{axis}
\onslide<2->{
\begin{axis}[onda, width=5.4cm, height=0.6cm, at={(0,-1.35cm)}, xmin=-6, xmax=6, ymin=0, ymax=1.3,
             every axis y label/.style={at={(rel axis cs:0,0.3)}, anchor=east, font=\scriptsize},
             xlabel={$f$}, ylabel={$S_n(f+f_c)$}]
  \addplot[cPurple, thick, fill=cPurple!20, samples=2] coordinates {(-1,0) (-1,0.5) (0,0.5) (0,0)};
  \addplot[cPurple, thick, fill=cPurple!20, samples=2] coordinates {(-5.5,0) (-5.5,0.5) (-4.5,0.5) (-4.5,0)};
  \node[font=\tiny, anchor=north] at (axis cs:-4.5,0) {$-2f_c$};
  \node[font=\tiny, anchor=north] at (axis cs:-1,0) {$-W$};
\end{axis}}
\onslide<3->{
\begin{axis}[onda, width=5.4cm, height=0.6cm, at={(0,-2.7cm)}, xmin=-6, xmax=6, ymin=0, ymax=1.3,
             every axis y label/.style={at={(rel axis cs:0,0.3)}, anchor=east, font=\scriptsize},
             xlabel={$f$}, ylabel={$S_x(f)$}]
  \addplot[cOrange, thick, fill=cOrange!20, samples=2] coordinates {(0,0) (0,0.5) (1,0.5) (1,0)};
  \addplot[cPurple, thick, fill=cPurple!20, samples=2] coordinates {(-1,0) (-1,0.5) (0,0.5) (0,0)};
  \addplot[cOrange, thick, fill=cOrange!20, samples=2] coordinates {(4.5,0) (4.5,0.5) (5.5,0.5) (5.5,0)};
  \addplot[cPurple, thick, fill=cPurple!20, samples=2] coordinates {(-5.5,0) (-5.5,0.5) (-4.5,0.5) (-4.5,0)};
  \only<4->{
    \draw[cGreen, thick, dashed] (axis cs:-6,0) -- (axis cs:-1,0) -- (axis cs:-1,1.2)
      -- (axis cs:1,1.2) -- (axis cs:1,0) -- (axis cs:6,0);
    \draw[cRed, very thick] (axis cs:4.4,-0.1) -- (axis cs:5.6,0.6) (axis cs:4.4,0.6) -- (axis cs:5.6,-0.1);
    \draw[cRed, very thick] (axis cs:-5.6,-0.1) -- (axis cs:-4.4,0.6) (axis cs:-5.6,0.6) -- (axis cs:-4.4,-0.1);
  }
\end{axis}}
\onslide<5->{
\begin{axis}[onda, width=5.4cm, height=0.6cm, at={(0,-4.05cm)}, xmin=-6, xmax=6, ymin=0, ymax=1.3,
             every axis y label/.style={at={(rel axis cs:0,0.3)}, anchor=east, font=\scriptsize},
             xlabel={$f$}, ylabel={$S_{n_c}(f)$}]
  \addplot[cRed, thick, fill=cRed!20, samples=2] coordinates {(-1,0) (-1,0.5) (1,0.5) (1,0)};
  \node[font=\tiny, anchor=north] at (axis cs:-1,0) {$-W$};
  \node[font=\tiny, anchor=north] at (axis cs:1,0) {$W$};
  \node[cRed, font=\tiny, anchor=west] at (axis cs:1.1,0.5) {$N_0/2$};
\end{axis}}
\end{tikzpicture}

\column{0.45\textwidth}
\small\compacto
Mesma regra do DSB-SC: $S_x(f) = S_n(f-f_c) + S_n(f+f_c)$.

\only<1>{A faixa de $[f_c, f_c+W]$ cai em $[0, W]$.}%
\only<2>{A faixa de $[-f_c-W, -f_c]$ cai em $[-W, 0]$.}%
\only<3>{As duas faixas não se empilham: encostam em $f = 0$. A altura fica $N_0/2$, metade do DSB-SC.}%
\only<4>{O passa-baixas guarda $|f|\le W$ e elimina as faixas de $\pm 2f_c$.}%
\only<5->{
\[
\overline{n_c^2} = \frac{N_0}{2}\cdot 2W = N_0W
\]
Mesma potência do ruído do ponto 1, como no DSB-SC: $\overline{n_c^2} = \overline{n^2}$.}
\end{columns}

\end{frame}

% ============================================

\begin{frame}{SSB, ponto 2: depois da detecção}

\small\compacto
Entrada do detector, agrupando por portadora:
\[
r(t) = \big[A_c m + n_c\big]\cos(2\pi f_c t) - \big[\tikzmarknode{qm}{A_c m_h + n_s}\big]\sin(2\pi f_c t)
\]
\par\vspace{0.5cm}%
\anotacoes{\only<2->{
  \node[nota=cOrange, anchor=north east] (a) at ([xshift=0.6cm,yshift=-0.25cm]qm.south) {vezes $\cos$ vira $\tfrac12\sin 4\pi f_c t$: o PB elimina};
  \draw[setinha, cOrange] ([xshift=-0.6cm]a.north east) -- (qm.south);
}}%
\onslide<3->{
\[
y(t) = \tfrac12 A_c\,m(t) + \tfrac12\,n_c(t)
\]}%
\onslide<4->{
\[
S_o = \frac{A_c^2\,\overline{m^2}}{4} = \frac{S_i}{4}, \qquad
N_o = \frac14\,\overline{n_c^2} = \frac{N_0W}{4}
\]}%
\onslide<5->{
\[
\left(\frac{S}{N}\right)_2 = \frac{S_i/4}{N_0W/4}
\only<6->{= \frac{S_i}{N_0W} = \gamma}
\]}%
\onslide<7->{Ganho de detecção 1: o detector descarta $n_s$, mas também descarta $m_h$, que levava metade de $S_i$.}

\end{frame}

% ============================================

\begin{frame}{DSB-SC e SSB lado a lado}

\small
\begin{center}
\begin{tabular}{l c c c c c}
\toprule
 & banda & ruído no ponto 1 & SNR no ponto 1 & SNR no ponto 2 & ganho de detecção \\
\midrule
banda base & $W$ & $N_0W$ & $\gamma$ & $\gamma$ & -- \\
\pause
DSB-SC & $2W$ & $2N_0W$ & $\gamma/2$ & $\gamma$ & 2 \\
\pause
SSB & $W$ & $N_0W$ & $\gamma$ & $\gamma$ & 1 \\
\bottomrule
\end{tabular}
\end{center}

\begin{itemize}
\item<4-> No DSB-SC, as duas bandas laterais carregam a mesma mensagem e se somam em fase no detector; o ruído das duas bandas é independente e se soma em potência.
\item<5-> No SSB não há essa soma, mas também não entra o ruído da outra banda.
\item<6-> As duas chegam a $\gamma$. O SSB usa metade da banda; o preço é o filtro e o modulador mais difíceis.
\end{itemize}

\end{frame}

% ============================================

% ============================================

\begin{frame}{Resumo: modulações lineares}

\small
\begin{center}
\begin{tabular}{l c c l}
\toprule
 & banda & detector & SNR de saída \\
\midrule
banda base & $W$ & -- & $\gamma$ \\
\pause
DSB-SC & $2W$ & coerente & $\gamma$ \\
\pause
AM convencional & $2W$ & envoltória & $\eta\gamma \le \gamma/3$, com limiar \\
\pause
SSB & $W$ & coerente & $\gamma$ \\
\bottomrule
\end{tabular}
\end{center}

\begin{itemize}
\item<5-> Com detecção coerente, nenhuma modulação linear supera a banda base: o melhor possível é $\gamma$.
\item<6-> O AM convencional paga a portadora para usar um receptor simples.
\item<7-> Para passar de $\gamma$ é preciso trocar banda por SNR, como faz o FM.
\end{itemize}

\end{frame}

% ============================================

\begin{frame}{Eficiência espectral dos esquemas}

\small\compacto
Banda da mensagem sobre banda transmitida:
\[
\rho = \frac{W}{B_T}
\]
\begin{center}
\begin{tabular}{l c c c}
\toprule
 & $B_T$ & $\rho$ & SNR$/\gamma$ \\
\midrule
\pause
banda base & $W$ & 1 & 1 \\
\pause
SSB & $W$ & 1 & 1 \\
\pause
VSB & $W + f_v$ & ${\approx}\,0{,}8$ & ${\approx}\,1$ \\
\pause
DSB-SC & $2W$ & $1/2$ & 1 \\
\pause
AM, $\mu = 1$ & $2W$ & $1/2$ & $1/3$ \\
\pause
FM, $\beta = 2$ & $6W$ & $1/6$ & 6 \\
FM, $\beta = 5$ & $12W$ & $1/12$ & $37{,}5$ \\
\bottomrule
\end{tabular}
\end{center}
\onslide<1->

\onslide<4->{VSB com vestígio $f_v \approx W/4$.}
\onslide<7->{FM com tom: $B_T = 2(\beta+1)W$ por Carson e SNR $= \tfrac32\beta^2\gamma$.}

\end{frame}

% ============================================

\begin{frame}{Banda contra SNR}

\begin{columns}[c]
\column{0.58\textwidth}
\centering
\begin{tikzpicture}
\begin{axis}[scale only axis, width=6cm, height=4.4cm,
             xmin=0, xmax=1.15, ymin=-7, ymax=18, axis lines=left,
             xlabel={$\rho = W/B_T$}, ylabel={SNR$/\gamma$ [dB]},
             xtick={0,0.25,0.5,0.75,1}, ytick={-5,0,5,10,15},
             tick label style={font=\tiny}, label style={font=\scriptsize},
             grid=major, grid style={gray!20}, clip=false]
\addplot[gray, dashed, domain=0:1.15] {0};
\addplot[only marks, mark=*, cBlue] coordinates {(1,0)}
  node[above, font=\scriptsize, cBlue] {banda base};
\only<2->{\addplot[only marks, mark=o, mark size=3.5pt, cBlue] coordinates {(1,0)}
  node[below, font=\scriptsize, cBlue] {SSB};}
\only<3->{\addplot[only marks, mark=*, cGreen] coordinates {(0.8,0)}
  node[above, font=\scriptsize, cGreen] {VSB};}
\only<4->{\addplot[only marks, mark=*, cBlue] coordinates {(0.5,0)}
  node[above, font=\scriptsize, cBlue] {DSB-SC};}
\only<5->{\addplot[only marks, mark=*, cRed] coordinates {(0.5,-4.77)}
  node[right, font=\scriptsize, cRed] {AM, $-4{,}8$ dB};}
\only<6->{\addplot[only marks, mark=*, cOrange] coordinates {(1/6,7.78) (1/12,15.74)};
  \node[right, font=\scriptsize, cOrange] at (axis cs:1/6,7.78) {FM, $\beta=2$};
  \node[right, font=\scriptsize, cOrange] at (axis cs:1/12,15.74) {FM, $\beta=5$};}
\end{axis}
\end{tikzpicture}

\column{0.42\textwidth}
\small
\onslide<7->{Nas modulações lineares, gastar mais banda não melhora a SNR: o ponto ótimo é o SSB, com $\rho = 1$ e SNR $= \gamma$.}

\vspace{0.3cm}
\onslide<8->{O AM paga a portadora: mesma banda do DSB-SC e SNR 3 vezes menor.}

\vspace{0.3cm}
\onslide<9->{O FM é o único que vai para a esquerda e para cima: troca eficiência espectral por SNR.}
\end{columns}

\end{frame}

% ============================================
